\subsubsection{Task 3 - Character Segmentation}
Character segmentation is the process of locating where individual characters begin and end within a row of text. Unlike row segmentation, which is a one-dimensional problem (finding row boundaries along the vertical axis), character segmentation is slightly more complex because characters in a row can have varying widths and can nearly touch one another. A clear example of this is the width difference between a narrow character like "I" and a wide character like "W”. \\

The function computes a vertical projection for every extracted row image. This projection corresponds to the sum of active pixels in each column, allowing us to distinguish between regions containing characters and blank spaces.\\\

The projection is then binarized by applying a threshold, producing a binary vector in which values equal to 1 correspond to columns containing foreground pixels, while values equal to 0 represent background regions. This additional binarization step reduces the influence of noise and facilitates the detection of character boundaries.\\

Character segmentation is performed by analyzing transitions in the binary projection using MATLAB’s \texttt{diff()} function. Specifically:

\begin{enumerate}[label=\alph*), leftmargin=1cm]
    \item \texttt{diff = 1}: indicates the beginning of a character region.
    \item \texttt{diff = -1}: indicates the end of a character region.
\end{enumerate}
From these transitions, the start and end positions of all character candidates are obtained.
